\section{Wpływ liczby jednostek wykonawczych na wydajność}
\label{sec:LiczbaRdzeni}

Kolejnym analizowanym aspektem był wpływ liczby wykorzystywanych jednostek
wykonawczych na czas działania równoległych wariantów badanych algorytmów.
Analizę przeprowadzono dla zbioru \texttt{random\_int} zawierającego 1~000~000 elementów, 
wykorzystując kolejno 2, 4, 8, 16 oraz 32 jednostki wykonawcze. 
Pozwoliło to określić, jak zmiana liczby jednostek wykonawczych wpływa na czas sortowania.

W celu porównania wpływu liczby jednostek wykonawczych na wydajność poszczególnych algorytmów 
zestawiono wyniki uzyskane dla wszystkich badanych konfiguracji.

\input{tables/chapter6/scalability_time}

Uzyskane wyniki wskazują wyraźne różnice pomiędzy badanymi algorytmami.
Wzrost liczby jednostek wykonawczych nie zawsze przekładał się na poprawę uzyskiwanego czasu sortowania, a najkrótsze
czasy osiągano przy różnej liczbie wykorzystywanych jednostek.

W~przypadku Quick Sort najkrótszy średni czas wykonania uzyskano już przy wykorzystaniu 2 jednostek wykonawczych i~wyniósł on około 1,80~s.
Zwiększenie ich liczby do 4 skutkowało wydłużeniem czasu do około 2,01~s, natomiast dla 8 jednostek osiągnął on około 2,22~s. 
Dalsze zwiększanie liczby jednostek do 16 i~32 nie powodowało już wyraźnych różnic w czasie wykonania, który utrzymywał się na poziomie około 2,22~s. 
Wyniki wskazują, że dla badanej implementacji Quick Sort zwiększenie liczby jednostek wykonawczych 
powyżej 2 nie przynosiło dalszych korzyści pod względem czasu wykonania.

Dla Merge Sort zwiększenie liczby jednostek wykonawczych z 2 do 4 pozwoliło skrócić czas z około 2,33~s do 2,15~s. 
Zwiększenie ich liczby do 8 nie przyniosło dalszej poprawy, a czas wyniósł około 2,33~s.
Dla 16 oraz 32 jednostek uzyskano odpowiednio około 2,34~s oraz 2,31~s.
Najkorzystniejszy wynik dla Merge Sort osiągnięto więc przy 4 jednostkach wykonawczych.

Odmienną tendencję zaobserwowano dla Bucket Sort. 
Zwiększanie liczby jednostek wykonawczych od 2 do 16 prowadziło do systematycznego skracania czasu sortowania. 
Średni czas skrócił się z około 1,88~s dla 2 jednostek do 1,63~s dla 4, 1,31~s dla 8 oraz 1,20~s dla 16 jednostek wykonawczych. 
Dla 32 jednostek nastąpiło natomiast wydłużenie czasu wykonania do około 1,48~s.
Najkorzystniejszy wynik dla Bucket Sort uzyskano zatem przy 16 jednostkach wykonawczych.

Sample Sort również początkowo charakteryzował się poprawą wydajności.
Średni czas wykonania zmniejszył się z około 3,96~s dla 2 jednostek do 2,89~s dla 4 oraz 2,37~s dla 8 jednostek wykonawczych. 
Przy 16 jednostkach uzyskano dalsze, niewielkie skrócenie czasu do około 2,44~s. 
Dla 32 jednostek nastąpiło natomiast wyraźne wydłużenie czasu wykonania do około 3,26~s.
Najlepszy wynik uzyskano zatem przy 8 jednostkach wykonawczych.

Zmiany czasu wykonania poszczególnych algorytmów wraz ze zwiększaniem liczby 
jednostek wykonawczych przedstawiono na rysunku~\ref{fig:time-cores-random-int}.

\begin{figure}[htbp]
    \centering
    \includegraphics[width=0.90\textwidth]
    {images/chapter6/execution_time/execution_time_vs_cores_random_int.png}
    \caption{Średni czas wykonania wariantów równoległych dla zbioru
    \texttt{random\_int} o rozmiarze 1~000~000 elementów w zależności od liczby jednostek wykonawczych. 
    Słupki błędów przedstawiają odchylenie standardowe. Źródło: opracowanie własne.}
    \label{fig:time-cores-random-int}
\end{figure}

Uzyskane wyniki wskazują, że najkorzystniejsza liczba jednostek wykonawczych różniła się w zależności od badanego algorytmu.
Quick Sort osiągnął najlepszy wynik przy 2 jednostkach, Merge Sort przy 4, 
Sample Sort przy 8, natomiast Bucket Sort przy 16 jednostkach wykonawczych.
Zwiększanie ich liczby powyżej tych wartości nie prowadziło do dalszej poprawy, 
a dla niektórych algorytmów wiązało się ze wzrostem czasu wykonania.

W~celu określenia, czy podobne zależności występują również dla innego typu danych, przeanalizowano wyniki dla zbioru \texttt{random\_float}.
Dla Merge Sort oraz Bucket Sort najkrótszy czas wykonania uzyskano przy tej samej 
liczbie jednostek wykonawczych co dla zbioru \texttt{random\_int}, odpowiednio przy 4 oraz 16 jednostkach. 
W~przypadku Quick Sort najlepszy wynik również uzyskano przy 2 jednostkach wykonawczych. 
Odmienną zależność zaobserwowano natomiast dla Sample Sort, 
dla którego najniższy średni czas wykonania przy danych zmiennoprzecinkowych uzyskano przy 16 jednostkach, 
podczas gdy dla danych całkowitych najlepszy wynik odnotowano przy 8 jednostkach. 
Różnica pomiędzy konfiguracjami 8 i~16 jednostek dla danych zmiennoprzecinkowych była jednak niewielka w~odniesieniu do~zmienności pomiarów.

Podobną tendencję zaobserwowano szczególnie dla Bucket Sort oraz Sample Sort.
Dla obu typów danych Bucket Sort skracał czas sortowania wraz ze wzrostem liczby 
jednostek wykonawczych aż do konfiguracji 16 jednostek, po której nastąpiło wydłużenie czasu wykonania.
W przypadku Sample Sort najkrótszy czas dla danych całkowitych uzyskiwał przy 8 jednostkach, a~dla danych zmiennoprzecinkowych przy 16 jednostkach. 
Zwiększenie liczby jednostek wykonawczych powyżej tych wartości prowadziło już do wydłużenia czasu wykonania.
Powtarzalność tych zależności dla dwóch typów danych sugeruje, że obserwowane zmiany wydajności mogły wynikać 
przede wszystkim ze sposobu realizacji równoległości badanych algorytmów, a~nie wyłącznie z~typu przetwarzanych wartości.

Uzyskane wyniki pokazują, że zwiększenie liczby wykorzystywanych jednostek nie prowadziło do proporcjonalnego wzrostu wydajności.
Dodatkowa diagnostyka wykonania przeprowadzona dla zbioru \texttt{random\_int} o~rozmiarze 1~000~000 elementów wykazała, 
że liczba zadeklarowanych jednostek wykonawczych nie odpowiadała bezpośrednio liczbie 
procesów roboczych faktycznie tworzonych przez badane algorytmy. 
W~przypadku Quick Sort utworzono 2 procesy potomne dla konfiguracji 2 jednostek,
6 procesów dla 4 jednostek oraz po 4 procesy dla konfiguracji 8, 16 i~32 jednostek wykonawczych. 
W~przypadku Merge Sort utworzono odpowiednio 2, 6, 2, 2 oraz 2 procesy potomne.

Uzyskane wyniki pokazują, że zwiększanie liczby zadeklarowanych jednostek wykonawczych 
nie zawsze prowadziło do~proporcjonalnego wzrostu liczby tworzonych procesów. 
W badanej implementacji parametr ten wyznaczał przede wszystkim maksymalną głębokość zrównoleglenia, 
natomiast rzeczywista liczba procesów zależała również od progów uruchamiania równoległego oraz rozmiaru powstałych fragmentów danych. 
W przypadku Merge Sort od konfiguracji 8 jednostek liczba utworzonych procesów pozostawała na poziomie dwóch. 
Dla Quick Sort również nie obserwowano dalszego wzrostu liczby procesów wraz ze zwiększaniem liczby zadeklarowanych jednostek od konfiguracji 8.

Dla porównania, w~przypadku Bucket Sort liczba utworzonych procesów potomnych odpowiadała liczbie zadeklarowanych jednostek wykonawczych. 
Sample Sort tworzył natomiast osobne zestawy procesów dla faz sortowania lokalnego oraz sortowania kubełków, 
przez co w~każdej z~tych faz liczba procesów odpowiadała liczbie zadeklarowanych jednostek wykonawczych. 
Oznacza to, że zależność pomiędzy liczbą zadeklarowanych jednostek wykonawczych a~liczbą faktycznie
tworzonych procesów była uzależniona od sposobu organizacji równoległości w~konkretnym algorytmie.

Przeprowadzona diagnostyka potwierdziła prawidłowe działanie algorytmów we wszystkich analizowanych konfiguracjach. 
Mniejsza od oczekiwanej liczba tworzonych procesów była więc skutkiem zastosowanego mechanizmu sterowania równoległością, 
a nie nieprawidłowego działania algorytmów.

Uzyskane wyniki pokazują, że zwiększenie liczby wykorzystywanych jednostek wykonawczych nie prowadziło do proporcjonalnego wzrostu wydajności.
Na osiągane rezultaty mogły wpływać również sposób zrównoleglenia konkretnego algorytmu oraz narzut związany z organizacją obliczeń równoległych. 
Wpływ tych czynników można ocenić również na podstawie wykorzystania zasobów oraz uzyskanego przyspieszenia i~efektywności.